\section{Hardware variation}
\label{si:hardware}

A trained network is a function of its data, its settings, its random seed and, on a GPU, the order in which
floating-point sums are accumulated. The aggregation over neighbouring atoms in each MEGNet block has no
deterministic GPU implementation, so the same run on two GPU models follows two slightly different training
trajectories, which can stop early at different passes and end in different weights. This section measures how
much that changes the predicted vacancy formation energies (extension; \texttt{docs/nrp/n1.md},
\texttt{results/nrp/nrp\_n1\_gpu\_repeat.json}).

The repeatability study (N1) reran the \numNrpNOneGpuRepeatCrossNodeSameClassSummaryN{} registered sweep runs of
resample \numCfgTuningAnchorResample{} at the two budgets \numCfgBudgetsHostsItemSix{} and
\numCfgBudgetsHostsItemFour{} (S and D, three seeds, originally registered network, unchanged settings) on four
GPU models, \numNrpNOneGpuRepeatNRuns{} runs in all. On one GPU model the runs are bitwise repeatable: a full
repeat on the same node reproduced every prediction
(\numNrpNOneGpuRepeatSameNodeRepeatSummaryBitwiseIdentical{} of \numNrpNOneGpuRepeatSameNodeRepeatSummaryN{} runs,
largest difference \numNrpNOneGpuRepeatSameNodeRepeatSummaryMaxAbsDiffEV{}~eV), and so did the same runs on a
second node with the same GPU model (\numNrpNOneGpuRepeatCrossNodeSameClassSummaryBitwiseIdentical{} of
\numNrpNOneGpuRepeatCrossNodeSameClassSummaryN{}). Between GPU models no run was bitwise identical, and single
predictions of one site differed by up to
\numNOneRecomputeCrossClassATenGpuFiveNrpMghpccOrgPassOneRtxAFourThousandRyGpuSixteenSdscOptiputerNetPassOneMaxAbsDiffEV{}~eV
between two classes, because a different stopping pass gives a different single network. The quantity that carries
the prediction in this work is the seed ensemble, and Table~\ref{tab:n1} gives its site MAE on each GPU model
minus that of the original sweep. At \numCfgBudgetsHostsItemSix{} training hosts the GPU model moves the ensemble
MAE of S by \numNrpNOneGpuRepeatEnsembleGeforceRtxThreeThousandNinetySSixHundredFiftyFourRyGpuTwoSdscOptiputerNetPassOneEnsembleMaeMinusRef{}
to \numNrpNOneGpuRepeatEnsembleLFourSSixHundredFiftyFourNautilusItGpuFiveFullertonEduPassOneEnsembleMaeMinusRef{}~eV
and that of D by \numNrpNOneGpuRepeatEnsembleGeforceRtxThreeThousandNinetyDStateSixHundredFiftyFourRyGpuTwoSdscOptiputerNetPassOneEnsembleMaeMinusRef{}
to \numNrpNOneGpuRepeatEnsembleRtxAFourThousandDStateSixHundredFiftyFourRyGpuSixteenSdscOptiputerNetPassOneEnsembleMaeMinusRef{}~eV;
at \numCfgBudgetsHostsItemFour{} hosts, where fewer DFT labels constrain the learned mapping, the shift of S
reaches \numNrpNOneGpuRepeatEnsembleRtxAFourThousandSTwoHundredRyGpuSixteenSdscOptiputerNetPassOneEnsembleMaeMinusRef{}~eV. The original sweep did
not record its GPU model, so its differences from the N1 runs mix GPU model with any difference in environment; the
comparisons within N1 share one code commit and one software environment.

\begin{table}[htbp]
\caption{GPU repeatability (N1; extension, descriptive). Seed-ensemble site MAE on each GPU model minus that of
the original sweep (eV), for the registered runs of resample \numCfgTuningAnchorResample{} of the originally
registered network; sweep MAE: the seed-ensemble MAE of the original sweep; seed SD: the standard deviation of the
single-run MAE over the three seeds of the original sweep, for scale. The \gpushortaten{} column holds the first pass on the
first node; the second pass and the second node are bitwise identical to it. No \gpufourty{} was free during the
study. Record: \texttt{results/nrp/nrp\_n1\_gpu\_repeat.json}, key \texttt{payload.ensemble}.}
\label{tab:n1}
\centering\small
% Generated by paper/si/tools/make_si_tables.py from results/nrp/nrp_n1_gpu_repeat.json (payload.ensemble). Do not edit.
\begin{tabular}{@{}llrrrrrr@{}}
\toprule
Model & Budget & \gpushortaten{} & \gpushortthreeoh{} & \gpushortafour{} & \gpushortlfour{} & Sweep MAE & Seed SD \\
\midrule
S & \numCfgBudgetsHostsItemSix{} & \numNrpNOneGpuRepeatEnsembleATenSSixHundredFiftyFourGpuFiveNrpMghpccOrgPassOneEnsembleMaeMinusRef{} & \numNrpNOneGpuRepeatEnsembleGeforceRtxThreeThousandNinetySSixHundredFiftyFourRyGpuTwoSdscOptiputerNetPassOneEnsembleMaeMinusRef{} & \numNrpNOneGpuRepeatEnsembleRtxAFourThousandSSixHundredFiftyFourRyGpuSixteenSdscOptiputerNetPassOneEnsembleMaeMinusRef{} & \numNrpNOneGpuRepeatEnsembleLFourSSixHundredFiftyFourNautilusItGpuFiveFullertonEduPassOneEnsembleMaeMinusRef{} & \numNrpNOneGpuRepeatEnsembleATenSSixHundredFiftyFourGpuFiveNrpMghpccOrgPassOneRefEnsembleMae{} & \numNrpNOneGpuRepeatEnsembleATenSSixHundredFiftyFourGpuFiveNrpMghpccOrgPassOneSeedSdResampleZero{} \\
S & \numCfgBudgetsHostsItemFour{} & \numNrpNOneGpuRepeatEnsembleATenSTwoHundredGpuFiveNrpMghpccOrgPassOneEnsembleMaeMinusRef{} & \numNrpNOneGpuRepeatEnsembleGeforceRtxThreeThousandNinetySTwoHundredRyGpuTwoSdscOptiputerNetPassOneEnsembleMaeMinusRef{} & \numNrpNOneGpuRepeatEnsembleRtxAFourThousandSTwoHundredRyGpuSixteenSdscOptiputerNetPassOneEnsembleMaeMinusRef{} & \numNrpNOneGpuRepeatEnsembleLFourSTwoHundredNautilusItGpuFiveFullertonEduPassOneEnsembleMaeMinusRef{} & \numNrpNOneGpuRepeatEnsembleATenSTwoHundredGpuFiveNrpMghpccOrgPassOneRefEnsembleMae{} & \numNrpNOneGpuRepeatEnsembleATenSTwoHundredGpuFiveNrpMghpccOrgPassOneSeedSdResampleZero{} \\
D & \numCfgBudgetsHostsItemSix{} & \numNrpNOneGpuRepeatEnsembleATenDStateSixHundredFiftyFourGpuFiveNrpMghpccOrgPassOneEnsembleMaeMinusRef{} & \numNrpNOneGpuRepeatEnsembleGeforceRtxThreeThousandNinetyDStateSixHundredFiftyFourRyGpuTwoSdscOptiputerNetPassOneEnsembleMaeMinusRef{} & \numNrpNOneGpuRepeatEnsembleRtxAFourThousandDStateSixHundredFiftyFourRyGpuSixteenSdscOptiputerNetPassOneEnsembleMaeMinusRef{} & \numNrpNOneGpuRepeatEnsembleLFourDStateSixHundredFiftyFourNautilusItGpuFiveFullertonEduPassOneEnsembleMaeMinusRef{} & \numNrpNOneGpuRepeatEnsembleATenDStateSixHundredFiftyFourGpuFiveNrpMghpccOrgPassOneRefEnsembleMae{} & \numNrpNOneGpuRepeatEnsembleATenDStateSixHundredFiftyFourGpuFiveNrpMghpccOrgPassOneSeedSdResampleZero{} \\
D & \numCfgBudgetsHostsItemFour{} & \numNrpNOneGpuRepeatEnsembleATenDStateTwoHundredGpuFiveNrpMghpccOrgPassOneEnsembleMaeMinusRef{} & \numNrpNOneGpuRepeatEnsembleGeforceRtxThreeThousandNinetyDStateTwoHundredRyGpuTwoSdscOptiputerNetPassOneEnsembleMaeMinusRef{} & \numNrpNOneGpuRepeatEnsembleRtxAFourThousandDStateTwoHundredRyGpuSixteenSdscOptiputerNetPassOneEnsembleMaeMinusRef{} & \numNrpNOneGpuRepeatEnsembleLFourDStateTwoHundredNautilusItGpuFiveFullertonEduPassOneEnsembleMaeMinusRef{} & \numNrpNOneGpuRepeatEnsembleATenDStateTwoHundredGpuFiveNrpMghpccOrgPassOneRefEnsembleMae{} & \numNrpNOneGpuRepeatEnsembleATenDStateTwoHundredGpuFiveNrpMghpccOrgPassOneSeedSdResampleZero{} \\
\bottomrule
\end{tabular}

\end{table}

The extensions trained on whichever GPU of the pool was free, so their runs mix GPU models. Figure~\ref{fig:gpu}
and Table~\ref{tab:gpu_class} give the advantage of the descriptors, $A = \mathrm{MAE}_\mathrm{S} -
\mathrm{MAE}_\mathrm{D}$, separately for the runs of each GPU model, using only cells in which the S run and the D
run of the same resample or fold, budget and seed ran on the same model (extension, descriptive, no test). On the
primary network (v2, post hoc) at \numCfgBudgetsHostsItemSix{} hosts the cluster classes give $A$ of
\numPrecisionABFourHundredBSixHundredFiftyFourDescriptiveGpuClassPrecisionSixHundredFiftyFourVTwoRTXFourThousandNinetyAPairedCellsA{}
to \numPrecisionABFourHundredBSixHundredFiftyFourDescriptiveGpuClassPrecisionSixHundredFiftyFourVTwoATenAPairedCellsA{}~eV,
against \numPrecisionABFourHundredBSixHundredFiftyFourDescriptiveGpuClassPrecisionSixHundredFiftyFourVTwoOriginalSweepAPairedCellsA{}~eV
for the original sweep runs; on the originally registered network the classes lie closer together. The GPU model
was not randomised with respect to resample, fold or seed, and the original sweep holds only the first resamples and
seeds, so a difference between classes cannot be separated from a difference in the resamples they ran and is not
read as an effect of the hardware. One cluster class, \gpufourty{}, is absent from N1 and ran many of the v2
precision runs. The non-inferiority statements of the article rest on the pooled runs at \numCfgBudgetsHostsItemSix{} training
hosts and at the cross-validated training sizes, on both networks; nothing in this section changes them, but every prediction carries a GPU
component of the size of Table~\ref{tab:n1} that is not separable from seed and resample variation.

\begin{figure}[htbp]
\centering
\sifig{fig_si_gpu_class}
\caption{The advantage $A$ by GPU class of the training runs (extension, descriptive, no test). $A =
\mathrm{MAE}_\mathrm{S} - \mathrm{MAE}_\mathrm{D}$ (eV) of single runs, averaged over the paired cells in which the S
run and the D run of the same resample or fold, budget and seed ran on the same GPU class. (\textbf{a}, \textbf{b})
Precision extension at the two largest budgets; (\textbf{c}, \textbf{d}) grouped $K$-fold cross-validation at the two
fold counts. The classes on the horizontal axis are the GPU models of the cluster pool, read from the run records,
and ``sweep'', the registered sweep runs of the original resamples and seeds (precision only). GPU class was not
randomised with respect to resample, fold or seed, so differences between classes are not separable from them and
are not read as hardware effects. Markers: mean over the paired cells; bars: descriptive percentile intervals of a
bootstrap over the cells; classes without a paired cell are not drawn. Filled markers and solid bars: primary network
(v2), adopted after diagnostics of the originally registered network; open markers and dashed bars: originally
registered network (v1). Diamonds: precision extension; triangles: grouped $K$-fold cross-validation. Dashed line:
the tolerance $\delta$; dotted line: zero. Values in Table~\ref{tab:gpu_class}.}
\label{fig:gpu}
\end{figure}

{\footnotesize\tighttabs
% Generated by paper/si/tools/make_si_tables.py from results/precision_a_b400_b654.json, results/cv_extended_budget.json (payload.descriptive.gpu_class). Do not edit.
\begin{longtable}{@{}lllllrrl@{}}
\caption{The advantage $A$ by GPU class of the training runs (extension, descriptive, no test). $A$ is the mean over paired cells (the S and D runs of the same resample or fold, budget and seed, both on that class) of the single-run difference MAE$_\mathrm{S}$ $-$ MAE$_\mathrm{D}$ (eV); the interval is a descriptive percentile interval of a bootstrap over the cells. Runs: all runs of S and of D on that class at that point. ``Original sweep'' holds the registered sweep runs of the original resamples and seeds. Records: \recprecision{} and \reccv{}, key \texttt{payload.descriptive.gpu\_class}.}\label{tab:gpu_class}\\
\toprule
Design & Point & Network & GPU class & Runs S / D & Cells & $A$ & Interval \\
\midrule
\endfirsthead
\toprule
Design & Point & Network & GPU class & Runs S / D & Cells & $A$ & Interval \\
\midrule
\endhead
\bottomrule
\endlastfoot
precision & \numCfgBudgetsHostsItemFive{} & v2 & \gpushortaten{} & \numPrecisionABFourHundredBSixHundredFiftyFourDescriptiveGpuClassPrecisionFourHundredVTwoATenNRunsS{} / \numPrecisionABFourHundredBSixHundredFiftyFourDescriptiveGpuClassPrecisionFourHundredVTwoATenNRunsD{} & \numPrecisionABFourHundredBSixHundredFiftyFourDescriptiveGpuClassPrecisionFourHundredVTwoATenNPairedCells{} & \numPrecisionABFourHundredBSixHundredFiftyFourDescriptiveGpuClassPrecisionFourHundredVTwoATenAPairedCellsA{} & \ci{\numPrecisionABFourHundredBSixHundredFiftyFourDescriptiveGpuClassPrecisionFourHundredVTwoATenAPairedCellsCiNinetyFiveDescriptiveItemOne{}}{\numPrecisionABFourHundredBSixHundredFiftyFourDescriptiveGpuClassPrecisionFourHundredVTwoATenAPairedCellsCiNinetyFiveDescriptiveItemTwo{}} \\
 &  & v2 & \gpushortthreeoh{} & \numPrecisionABFourHundredBSixHundredFiftyFourDescriptiveGpuClassPrecisionFourHundredVTwoRTXThreeThousandNinetyNRunsS{} / \numPrecisionABFourHundredBSixHundredFiftyFourDescriptiveGpuClassPrecisionFourHundredVTwoRTXThreeThousandNinetyNRunsD{} & \numPrecisionABFourHundredBSixHundredFiftyFourDescriptiveGpuClassPrecisionFourHundredVTwoRTXThreeThousandNinetyNPairedCells{} & \numPrecisionABFourHundredBSixHundredFiftyFourDescriptiveGpuClassPrecisionFourHundredVTwoRTXThreeThousandNinetyAPairedCellsA{} & \ci{\numPrecisionABFourHundredBSixHundredFiftyFourDescriptiveGpuClassPrecisionFourHundredVTwoRTXThreeThousandNinetyAPairedCellsCiNinetyFiveDescriptiveItemOne{}}{\numPrecisionABFourHundredBSixHundredFiftyFourDescriptiveGpuClassPrecisionFourHundredVTwoRTXThreeThousandNinetyAPairedCellsCiNinetyFiveDescriptiveItemTwo{}} \\
 &  & v2 & \gpushortlfour{} & \numPrecisionABFourHundredBSixHundredFiftyFourDescriptiveGpuClassPrecisionFourHundredVTwoLFourNRunsS{} / \numPrecisionABFourHundredBSixHundredFiftyFourDescriptiveGpuClassPrecisionFourHundredVTwoLFourNRunsD{} & \numPrecisionABFourHundredBSixHundredFiftyFourDescriptiveGpuClassPrecisionFourHundredVTwoLFourNPairedCells{} & \numPrecisionABFourHundredBSixHundredFiftyFourDescriptiveGpuClassPrecisionFourHundredVTwoLFourAPairedCellsA{} & \ci{\numPrecisionABFourHundredBSixHundredFiftyFourDescriptiveGpuClassPrecisionFourHundredVTwoLFourAPairedCellsCiNinetyFiveDescriptiveItemOne{}}{\numPrecisionABFourHundredBSixHundredFiftyFourDescriptiveGpuClassPrecisionFourHundredVTwoLFourAPairedCellsCiNinetyFiveDescriptiveItemTwo{}} \\
 &  & v2 & \gpushortfourty{} (not in N1) & \numPrecisionABFourHundredBSixHundredFiftyFourDescriptiveGpuClassPrecisionFourHundredVTwoRTXFourThousandNinetyNRunsS{} / \numPrecisionABFourHundredBSixHundredFiftyFourDescriptiveGpuClassPrecisionFourHundredVTwoRTXFourThousandNinetyNRunsD{} & \numPrecisionABFourHundredBSixHundredFiftyFourDescriptiveGpuClassPrecisionFourHundredVTwoRTXFourThousandNinetyNPairedCells{} & \numPrecisionABFourHundredBSixHundredFiftyFourDescriptiveGpuClassPrecisionFourHundredVTwoRTXFourThousandNinetyAPairedCellsA{} & \ci{\numPrecisionABFourHundredBSixHundredFiftyFourDescriptiveGpuClassPrecisionFourHundredVTwoRTXFourThousandNinetyAPairedCellsCiNinetyFiveDescriptiveItemOne{}}{\numPrecisionABFourHundredBSixHundredFiftyFourDescriptiveGpuClassPrecisionFourHundredVTwoRTXFourThousandNinetyAPairedCellsCiNinetyFiveDescriptiveItemTwo{}} \\
 &  & v2 & original sweep & \numPrecisionABFourHundredBSixHundredFiftyFourDescriptiveGpuClassPrecisionFourHundredVTwoOriginalSweepNRunsS{} / \numPrecisionABFourHundredBSixHundredFiftyFourDescriptiveGpuClassPrecisionFourHundredVTwoOriginalSweepNRunsD{} & \numPrecisionABFourHundredBSixHundredFiftyFourDescriptiveGpuClassPrecisionFourHundredVTwoOriginalSweepNPairedCells{} & \numPrecisionABFourHundredBSixHundredFiftyFourDescriptiveGpuClassPrecisionFourHundredVTwoOriginalSweepAPairedCellsA{} & \ci{\numPrecisionABFourHundredBSixHundredFiftyFourDescriptiveGpuClassPrecisionFourHundredVTwoOriginalSweepAPairedCellsCiNinetyFiveDescriptiveItemOne{}}{\numPrecisionABFourHundredBSixHundredFiftyFourDescriptiveGpuClassPrecisionFourHundredVTwoOriginalSweepAPairedCellsCiNinetyFiveDescriptiveItemTwo{}} \\
 &  & v1 & \gpushortaten{} & \numPrecisionABFourHundredBSixHundredFiftyFourDescriptiveGpuClassPrecisionFourHundredVOneATenNRunsS{} / \numPrecisionABFourHundredBSixHundredFiftyFourDescriptiveGpuClassPrecisionFourHundredVOneATenNRunsD{} & \numPrecisionABFourHundredBSixHundredFiftyFourDescriptiveGpuClassPrecisionFourHundredVOneATenNPairedCells{} & \numPrecisionABFourHundredBSixHundredFiftyFourDescriptiveGpuClassPrecisionFourHundredVOneATenAPairedCellsA{} & \ci{\numPrecisionABFourHundredBSixHundredFiftyFourDescriptiveGpuClassPrecisionFourHundredVOneATenAPairedCellsCiNinetyFiveDescriptiveItemOne{}}{\numPrecisionABFourHundredBSixHundredFiftyFourDescriptiveGpuClassPrecisionFourHundredVOneATenAPairedCellsCiNinetyFiveDescriptiveItemTwo{}} \\
 &  & v1 & \gpushortthreeoh{} & \numPrecisionABFourHundredBSixHundredFiftyFourDescriptiveGpuClassPrecisionFourHundredVOneRTXThreeThousandNinetyNRunsS{} / \numPrecisionABFourHundredBSixHundredFiftyFourDescriptiveGpuClassPrecisionFourHundredVOneRTXThreeThousandNinetyNRunsD{} & \numPrecisionABFourHundredBSixHundredFiftyFourDescriptiveGpuClassPrecisionFourHundredVOneRTXThreeThousandNinetyNPairedCells{} & \numPrecisionABFourHundredBSixHundredFiftyFourDescriptiveGpuClassPrecisionFourHundredVOneRTXThreeThousandNinetyAPairedCellsA{} & \ci{\numPrecisionABFourHundredBSixHundredFiftyFourDescriptiveGpuClassPrecisionFourHundredVOneRTXThreeThousandNinetyAPairedCellsCiNinetyFiveDescriptiveItemOne{}}{\numPrecisionABFourHundredBSixHundredFiftyFourDescriptiveGpuClassPrecisionFourHundredVOneRTXThreeThousandNinetyAPairedCellsCiNinetyFiveDescriptiveItemTwo{}} \\
 &  & v1 & \gpushortlfour{} & \numPrecisionABFourHundredBSixHundredFiftyFourDescriptiveGpuClassPrecisionFourHundredVOneLFourNRunsS{} / \numPrecisionABFourHundredBSixHundredFiftyFourDescriptiveGpuClassPrecisionFourHundredVOneLFourNRunsD{} & \numPrecisionABFourHundredBSixHundredFiftyFourDescriptiveGpuClassPrecisionFourHundredVOneLFourNPairedCells{} & \numPrecisionABFourHundredBSixHundredFiftyFourDescriptiveGpuClassPrecisionFourHundredVOneLFourAPairedCellsA{} & \ci{\numPrecisionABFourHundredBSixHundredFiftyFourDescriptiveGpuClassPrecisionFourHundredVOneLFourAPairedCellsCiNinetyFiveDescriptiveItemOne{}}{\numPrecisionABFourHundredBSixHundredFiftyFourDescriptiveGpuClassPrecisionFourHundredVOneLFourAPairedCellsCiNinetyFiveDescriptiveItemTwo{}} \\
 &  & v1 & \gpushortfourty{} (not in N1) & \numPrecisionABFourHundredBSixHundredFiftyFourDescriptiveGpuClassPrecisionFourHundredVOneRTXFourThousandNinetyNRunsS{} / \numPrecisionABFourHundredBSixHundredFiftyFourDescriptiveGpuClassPrecisionFourHundredVOneRTXFourThousandNinetyNRunsD{} & \numPrecisionABFourHundredBSixHundredFiftyFourDescriptiveGpuClassPrecisionFourHundredVOneRTXFourThousandNinetyNPairedCells{} & -- &  \\
 &  & v1 & original sweep & \numPrecisionABFourHundredBSixHundredFiftyFourDescriptiveGpuClassPrecisionFourHundredVOneOriginalSweepNRunsS{} / \numPrecisionABFourHundredBSixHundredFiftyFourDescriptiveGpuClassPrecisionFourHundredVOneOriginalSweepNRunsD{} & \numPrecisionABFourHundredBSixHundredFiftyFourDescriptiveGpuClassPrecisionFourHundredVOneOriginalSweepNPairedCells{} & \numPrecisionABFourHundredBSixHundredFiftyFourDescriptiveGpuClassPrecisionFourHundredVOneOriginalSweepAPairedCellsA{} & \ci{\numPrecisionABFourHundredBSixHundredFiftyFourDescriptiveGpuClassPrecisionFourHundredVOneOriginalSweepAPairedCellsCiNinetyFiveDescriptiveItemOne{}}{\numPrecisionABFourHundredBSixHundredFiftyFourDescriptiveGpuClassPrecisionFourHundredVOneOriginalSweepAPairedCellsCiNinetyFiveDescriptiveItemTwo{}} \\
\addlinespace
precision & \numCfgBudgetsHostsItemSix{} & v2 & \gpushortaten{} & \numPrecisionABFourHundredBSixHundredFiftyFourDescriptiveGpuClassPrecisionSixHundredFiftyFourVTwoATenNRunsS{} / \numPrecisionABFourHundredBSixHundredFiftyFourDescriptiveGpuClassPrecisionSixHundredFiftyFourVTwoATenNRunsD{} & \numPrecisionABFourHundredBSixHundredFiftyFourDescriptiveGpuClassPrecisionSixHundredFiftyFourVTwoATenNPairedCells{} & \numPrecisionABFourHundredBSixHundredFiftyFourDescriptiveGpuClassPrecisionSixHundredFiftyFourVTwoATenAPairedCellsA{} & \ci{\numPrecisionABFourHundredBSixHundredFiftyFourDescriptiveGpuClassPrecisionSixHundredFiftyFourVTwoATenAPairedCellsCiNinetyFiveDescriptiveItemOne{}}{\numPrecisionABFourHundredBSixHundredFiftyFourDescriptiveGpuClassPrecisionSixHundredFiftyFourVTwoATenAPairedCellsCiNinetyFiveDescriptiveItemTwo{}} \\
 &  & v2 & \gpushortthreeoh{} & \numPrecisionABFourHundredBSixHundredFiftyFourDescriptiveGpuClassPrecisionSixHundredFiftyFourVTwoRTXThreeThousandNinetyNRunsS{} / \numPrecisionABFourHundredBSixHundredFiftyFourDescriptiveGpuClassPrecisionSixHundredFiftyFourVTwoRTXThreeThousandNinetyNRunsD{} & \numPrecisionABFourHundredBSixHundredFiftyFourDescriptiveGpuClassPrecisionSixHundredFiftyFourVTwoRTXThreeThousandNinetyNPairedCells{} & \numPrecisionABFourHundredBSixHundredFiftyFourDescriptiveGpuClassPrecisionSixHundredFiftyFourVTwoRTXThreeThousandNinetyAPairedCellsA{} & \ci{\numPrecisionABFourHundredBSixHundredFiftyFourDescriptiveGpuClassPrecisionSixHundredFiftyFourVTwoRTXThreeThousandNinetyAPairedCellsCiNinetyFiveDescriptiveItemOne{}}{\numPrecisionABFourHundredBSixHundredFiftyFourDescriptiveGpuClassPrecisionSixHundredFiftyFourVTwoRTXThreeThousandNinetyAPairedCellsCiNinetyFiveDescriptiveItemTwo{}} \\
 &  & v2 & \gpushortlfour{} & \numPrecisionABFourHundredBSixHundredFiftyFourDescriptiveGpuClassPrecisionSixHundredFiftyFourVTwoLFourNRunsS{} / \numPrecisionABFourHundredBSixHundredFiftyFourDescriptiveGpuClassPrecisionSixHundredFiftyFourVTwoLFourNRunsD{} & \numPrecisionABFourHundredBSixHundredFiftyFourDescriptiveGpuClassPrecisionSixHundredFiftyFourVTwoLFourNPairedCells{} & -- &  \\
 &  & v2 & \gpushortfourty{} (not in N1) & \numPrecisionABFourHundredBSixHundredFiftyFourDescriptiveGpuClassPrecisionSixHundredFiftyFourVTwoRTXFourThousandNinetyNRunsS{} / \numPrecisionABFourHundredBSixHundredFiftyFourDescriptiveGpuClassPrecisionSixHundredFiftyFourVTwoRTXFourThousandNinetyNRunsD{} & \numPrecisionABFourHundredBSixHundredFiftyFourDescriptiveGpuClassPrecisionSixHundredFiftyFourVTwoRTXFourThousandNinetyNPairedCells{} & \numPrecisionABFourHundredBSixHundredFiftyFourDescriptiveGpuClassPrecisionSixHundredFiftyFourVTwoRTXFourThousandNinetyAPairedCellsA{} & \ci{\numPrecisionABFourHundredBSixHundredFiftyFourDescriptiveGpuClassPrecisionSixHundredFiftyFourVTwoRTXFourThousandNinetyAPairedCellsCiNinetyFiveDescriptiveItemOne{}}{\numPrecisionABFourHundredBSixHundredFiftyFourDescriptiveGpuClassPrecisionSixHundredFiftyFourVTwoRTXFourThousandNinetyAPairedCellsCiNinetyFiveDescriptiveItemTwo{}} \\
 &  & v2 & original sweep & \numPrecisionABFourHundredBSixHundredFiftyFourDescriptiveGpuClassPrecisionSixHundredFiftyFourVTwoOriginalSweepNRunsS{} / \numPrecisionABFourHundredBSixHundredFiftyFourDescriptiveGpuClassPrecisionSixHundredFiftyFourVTwoOriginalSweepNRunsD{} & \numPrecisionABFourHundredBSixHundredFiftyFourDescriptiveGpuClassPrecisionSixHundredFiftyFourVTwoOriginalSweepNPairedCells{} & \numPrecisionABFourHundredBSixHundredFiftyFourDescriptiveGpuClassPrecisionSixHundredFiftyFourVTwoOriginalSweepAPairedCellsA{} & \ci{\numPrecisionABFourHundredBSixHundredFiftyFourDescriptiveGpuClassPrecisionSixHundredFiftyFourVTwoOriginalSweepAPairedCellsCiNinetyFiveDescriptiveItemOne{}}{\numPrecisionABFourHundredBSixHundredFiftyFourDescriptiveGpuClassPrecisionSixHundredFiftyFourVTwoOriginalSweepAPairedCellsCiNinetyFiveDescriptiveItemTwo{}} \\
 &  & v1 & \gpushortaten{} & \numPrecisionABFourHundredBSixHundredFiftyFourDescriptiveGpuClassPrecisionSixHundredFiftyFourVOneATenNRunsS{} / \numPrecisionABFourHundredBSixHundredFiftyFourDescriptiveGpuClassPrecisionSixHundredFiftyFourVOneATenNRunsD{} & \numPrecisionABFourHundredBSixHundredFiftyFourDescriptiveGpuClassPrecisionSixHundredFiftyFourVOneATenNPairedCells{} & \numPrecisionABFourHundredBSixHundredFiftyFourDescriptiveGpuClassPrecisionSixHundredFiftyFourVOneATenAPairedCellsA{} & \ci{\numPrecisionABFourHundredBSixHundredFiftyFourDescriptiveGpuClassPrecisionSixHundredFiftyFourVOneATenAPairedCellsCiNinetyFiveDescriptiveItemOne{}}{\numPrecisionABFourHundredBSixHundredFiftyFourDescriptiveGpuClassPrecisionSixHundredFiftyFourVOneATenAPairedCellsCiNinetyFiveDescriptiveItemTwo{}} \\
 &  & v1 & \gpushortthreeoh{} & \numPrecisionABFourHundredBSixHundredFiftyFourDescriptiveGpuClassPrecisionSixHundredFiftyFourVOneRTXThreeThousandNinetyNRunsS{} / \numPrecisionABFourHundredBSixHundredFiftyFourDescriptiveGpuClassPrecisionSixHundredFiftyFourVOneRTXThreeThousandNinetyNRunsD{} & \numPrecisionABFourHundredBSixHundredFiftyFourDescriptiveGpuClassPrecisionSixHundredFiftyFourVOneRTXThreeThousandNinetyNPairedCells{} & \numPrecisionABFourHundredBSixHundredFiftyFourDescriptiveGpuClassPrecisionSixHundredFiftyFourVOneRTXThreeThousandNinetyAPairedCellsA{} & \ci{\numPrecisionABFourHundredBSixHundredFiftyFourDescriptiveGpuClassPrecisionSixHundredFiftyFourVOneRTXThreeThousandNinetyAPairedCellsCiNinetyFiveDescriptiveItemOne{}}{\numPrecisionABFourHundredBSixHundredFiftyFourDescriptiveGpuClassPrecisionSixHundredFiftyFourVOneRTXThreeThousandNinetyAPairedCellsCiNinetyFiveDescriptiveItemTwo{}} \\
 &  & v1 & \gpushortlfour{} & \numPrecisionABFourHundredBSixHundredFiftyFourDescriptiveGpuClassPrecisionSixHundredFiftyFourVOneLFourNRunsS{} / \numPrecisionABFourHundredBSixHundredFiftyFourDescriptiveGpuClassPrecisionSixHundredFiftyFourVOneLFourNRunsD{} & \numPrecisionABFourHundredBSixHundredFiftyFourDescriptiveGpuClassPrecisionSixHundredFiftyFourVOneLFourNPairedCells{} & \numPrecisionABFourHundredBSixHundredFiftyFourDescriptiveGpuClassPrecisionSixHundredFiftyFourVOneLFourAPairedCellsA{} & \ci{\numPrecisionABFourHundredBSixHundredFiftyFourDescriptiveGpuClassPrecisionSixHundredFiftyFourVOneLFourAPairedCellsCiNinetyFiveDescriptiveItemOne{}}{\numPrecisionABFourHundredBSixHundredFiftyFourDescriptiveGpuClassPrecisionSixHundredFiftyFourVOneLFourAPairedCellsCiNinetyFiveDescriptiveItemTwo{}} \\
 &  & v1 & \gpushortfourty{} (not in N1) & \numPrecisionABFourHundredBSixHundredFiftyFourDescriptiveGpuClassPrecisionSixHundredFiftyFourVOneRTXFourThousandNinetyNRunsS{} / \numPrecisionABFourHundredBSixHundredFiftyFourDescriptiveGpuClassPrecisionSixHundredFiftyFourVOneRTXFourThousandNinetyNRunsD{} & \numPrecisionABFourHundredBSixHundredFiftyFourDescriptiveGpuClassPrecisionSixHundredFiftyFourVOneRTXFourThousandNinetyNPairedCells{} & \numPrecisionABFourHundredBSixHundredFiftyFourDescriptiveGpuClassPrecisionSixHundredFiftyFourVOneRTXFourThousandNinetyAPairedCellsA{} & \ci{\numPrecisionABFourHundredBSixHundredFiftyFourDescriptiveGpuClassPrecisionSixHundredFiftyFourVOneRTXFourThousandNinetyAPairedCellsCiNinetyFiveDescriptiveItemOne{}}{\numPrecisionABFourHundredBSixHundredFiftyFourDescriptiveGpuClassPrecisionSixHundredFiftyFourVOneRTXFourThousandNinetyAPairedCellsCiNinetyFiveDescriptiveItemTwo{}} \\
 &  & v1 & original sweep & \numPrecisionABFourHundredBSixHundredFiftyFourDescriptiveGpuClassPrecisionSixHundredFiftyFourVOneOriginalSweepNRunsS{} / \numPrecisionABFourHundredBSixHundredFiftyFourDescriptiveGpuClassPrecisionSixHundredFiftyFourVOneOriginalSweepNRunsD{} & \numPrecisionABFourHundredBSixHundredFiftyFourDescriptiveGpuClassPrecisionSixHundredFiftyFourVOneOriginalSweepNPairedCells{} & \numPrecisionABFourHundredBSixHundredFiftyFourDescriptiveGpuClassPrecisionSixHundredFiftyFourVOneOriginalSweepAPairedCellsA{} & \ci{\numPrecisionABFourHundredBSixHundredFiftyFourDescriptiveGpuClassPrecisionSixHundredFiftyFourVOneOriginalSweepAPairedCellsCiNinetyFiveDescriptiveItemOne{}}{\numPrecisionABFourHundredBSixHundredFiftyFourDescriptiveGpuClassPrecisionSixHundredFiftyFourVOneOriginalSweepAPairedCellsCiNinetyFiveDescriptiveItemTwo{}} \\
\addlinespace
grouped CV & $K = \numCvExtendedBudgetDescriptiveGpuClassCvTenVTwoATenPoint{}$ & v2 & \gpushortaten{} & \numCvExtendedBudgetDescriptiveGpuClassCvTenVTwoATenNRunsS{} / \numCvExtendedBudgetDescriptiveGpuClassCvTenVTwoATenNRunsD{} & \numCvExtendedBudgetDescriptiveGpuClassCvTenVTwoATenNPairedCells{} & \numCvExtendedBudgetDescriptiveGpuClassCvTenVTwoATenAPairedCellsA{} & \ci{\numCvExtendedBudgetDescriptiveGpuClassCvTenVTwoATenAPairedCellsCiNinetyFiveDescriptiveItemOne{}}{\numCvExtendedBudgetDescriptiveGpuClassCvTenVTwoATenAPairedCellsCiNinetyFiveDescriptiveItemTwo{}} \\
 &  & v2 & \gpushortthreeoh{} & \numCvExtendedBudgetDescriptiveGpuClassCvTenVTwoRTXThreeThousandNinetyNRunsS{} / \numCvExtendedBudgetDescriptiveGpuClassCvTenVTwoRTXThreeThousandNinetyNRunsD{} & \numCvExtendedBudgetDescriptiveGpuClassCvTenVTwoRTXThreeThousandNinetyNPairedCells{} & \numCvExtendedBudgetDescriptiveGpuClassCvTenVTwoRTXThreeThousandNinetyAPairedCellsA{} & \ci{\numCvExtendedBudgetDescriptiveGpuClassCvTenVTwoRTXThreeThousandNinetyAPairedCellsCiNinetyFiveDescriptiveItemOne{}}{\numCvExtendedBudgetDescriptiveGpuClassCvTenVTwoRTXThreeThousandNinetyAPairedCellsCiNinetyFiveDescriptiveItemTwo{}} \\
 &  & v1 & \gpushortaten{} & \numCvExtendedBudgetDescriptiveGpuClassCvTenVOneATenNRunsS{} / \numCvExtendedBudgetDescriptiveGpuClassCvTenVOneATenNRunsD{} & \numCvExtendedBudgetDescriptiveGpuClassCvTenVOneATenNPairedCells{} & \numCvExtendedBudgetDescriptiveGpuClassCvTenVOneATenAPairedCellsA{} & \ci{\numCvExtendedBudgetDescriptiveGpuClassCvTenVOneATenAPairedCellsCiNinetyFiveDescriptiveItemOne{}}{\numCvExtendedBudgetDescriptiveGpuClassCvTenVOneATenAPairedCellsCiNinetyFiveDescriptiveItemTwo{}} \\
 &  & v1 & \gpushortlfour{} & \numCvExtendedBudgetDescriptiveGpuClassCvTenVOneLFourNRunsS{} / \numCvExtendedBudgetDescriptiveGpuClassCvTenVOneLFourNRunsD{} & \numCvExtendedBudgetDescriptiveGpuClassCvTenVOneLFourNPairedCells{} & \numCvExtendedBudgetDescriptiveGpuClassCvTenVOneLFourAPairedCellsA{} & \ci{\numCvExtendedBudgetDescriptiveGpuClassCvTenVOneLFourAPairedCellsCiNinetyFiveDescriptiveItemOne{}}{\numCvExtendedBudgetDescriptiveGpuClassCvTenVOneLFourAPairedCellsCiNinetyFiveDescriptiveItemTwo{}} \\
\addlinespace
grouped CV & $K = \numCvExtendedBudgetDescriptiveGpuClassCvTwentyVTwoATenPoint{}$ & v2 & \gpushortaten{} & \numCvExtendedBudgetDescriptiveGpuClassCvTwentyVTwoATenNRunsS{} / \numCvExtendedBudgetDescriptiveGpuClassCvTwentyVTwoATenNRunsD{} & \numCvExtendedBudgetDescriptiveGpuClassCvTwentyVTwoATenNPairedCells{} & \numCvExtendedBudgetDescriptiveGpuClassCvTwentyVTwoATenAPairedCellsA{} & \ci{\numCvExtendedBudgetDescriptiveGpuClassCvTwentyVTwoATenAPairedCellsCiNinetyFiveDescriptiveItemOne{}}{\numCvExtendedBudgetDescriptiveGpuClassCvTwentyVTwoATenAPairedCellsCiNinetyFiveDescriptiveItemTwo{}} \\
 &  & v2 & \gpushortthreeoh{} & \numCvExtendedBudgetDescriptiveGpuClassCvTwentyVTwoRTXThreeThousandNinetyNRunsS{} / \numCvExtendedBudgetDescriptiveGpuClassCvTwentyVTwoRTXThreeThousandNinetyNRunsD{} & \numCvExtendedBudgetDescriptiveGpuClassCvTwentyVTwoRTXThreeThousandNinetyNPairedCells{} & \numCvExtendedBudgetDescriptiveGpuClassCvTwentyVTwoRTXThreeThousandNinetyAPairedCellsA{} & \ci{\numCvExtendedBudgetDescriptiveGpuClassCvTwentyVTwoRTXThreeThousandNinetyAPairedCellsCiNinetyFiveDescriptiveItemOne{}}{\numCvExtendedBudgetDescriptiveGpuClassCvTwentyVTwoRTXThreeThousandNinetyAPairedCellsCiNinetyFiveDescriptiveItemTwo{}} \\
 &  & v1 & \gpushortaten{} & \numCvExtendedBudgetDescriptiveGpuClassCvTwentyVOneATenNRunsS{} / \numCvExtendedBudgetDescriptiveGpuClassCvTwentyVOneATenNRunsD{} & \numCvExtendedBudgetDescriptiveGpuClassCvTwentyVOneATenNPairedCells{} & \numCvExtendedBudgetDescriptiveGpuClassCvTwentyVOneATenAPairedCellsA{} & \ci{\numCvExtendedBudgetDescriptiveGpuClassCvTwentyVOneATenAPairedCellsCiNinetyFiveDescriptiveItemOne{}}{\numCvExtendedBudgetDescriptiveGpuClassCvTwentyVOneATenAPairedCellsCiNinetyFiveDescriptiveItemTwo{}} \\
 &  & v1 & \gpushortlfour{} & \numCvExtendedBudgetDescriptiveGpuClassCvTwentyVOneLFourNRunsS{} / \numCvExtendedBudgetDescriptiveGpuClassCvTwentyVOneLFourNRunsD{} & \numCvExtendedBudgetDescriptiveGpuClassCvTwentyVOneLFourNPairedCells{} & \numCvExtendedBudgetDescriptiveGpuClassCvTwentyVOneLFourAPairedCellsA{} & \ci{\numCvExtendedBudgetDescriptiveGpuClassCvTwentyVOneLFourAPairedCellsCiNinetyFiveDescriptiveItemOne{}}{\numCvExtendedBudgetDescriptiveGpuClassCvTwentyVOneLFourAPairedCellsCiNinetyFiveDescriptiveItemTwo{}} \\
\end{longtable}
}
\FloatBarrier
